% Minimal preamble for P2 cards taken from sources/comparability-partitions-2408.13523v2/source0.tex
% (arXiv:2408.13523v2, "Partitioning perfect graphs into comparability graphs").
% Every line below is copied mechanically from that file; nothing was invented, renamed or redefined.
%   line 2    : \documentclass[12pt]{article}
%   line 3    : \usepackage{amssymb,latexsym}
%   line 4    : \usepackage{amsmath}
%   line 5    : \usepackage[dvipdfm]{graphicx}  -- the dvipdfm driver option is dropped, because the
%               assembler forces the xelatex engine and a forced driver option contradicts it.
%   line 16   : \usepackage{amsfonts}
%   line 17   : \usepackage{amsthm}
%   lines 20-30: the author's theorem declarations
%   lines 32-41: the author's own macro definitions
% Omitted source preamble lines and why they are not needed for these body spans:
%   line 1 (provenance comment), lines 6-11 (tikz and its libraries: no drawing in the retained spans),
%   lines 7 and 9 (verbatim), lines 12-13 (page geometry), line 14 (rawfonts),
%   line 15 (psfrag), line 18 (hyperref with the pagebackref option: the retained spans print no
%   citation anchors and the assembly adds its own hyperref machinery), lines 43-61 (the author's
%   backref neutralisation and his margin-note macros, which depend on the omitted hyperref).
% The retained spans of the body use only \em, \bf, \cal, amsthm environments, amsmath displays and
% the macros defined in lines 32-41 below.
\documentclass[12pt]{article}
\usepackage{amssymb,latexsym}
\usepackage{amsmath}
\usepackage{graphicx}
\usepackage{amsfonts}
\usepackage{amsthm}

\newtheorem{theorem}{Theorem}%[section]
\newtheorem{problem}[theorem]{Problem}%[section]
\newtheorem{observation}[theorem]{Observation}
\newtheorem{lemma}[theorem]{Lemma}
\newtheorem{corollary}[theorem]{Corollary}
\newtheorem{definition}{Definition}
\newtheorem{example}[theorem]{Example}
\newtheorem{conjecture}{Conjecture}%[section]
\newtheorem{proposition}[theorem]{Proposition}
\newtheorem{claim}[theorem]{Claim}
\newtheorem{remark}[theorem]{Remark}

\def\qed{\ifhmode\unskip\nobreak\fi\quad\ifmmode\Box\else$\Box$\fi}
\def\pont{\hspace{-6pt}{\bf.\ }}
\def\qed{\hfill $\Box$}

\def\R{{\mbox{\boldmath$\mathcal R$}}}
% the author's second definition of \R, left commented out exactly as in the source
\def\F{{\mathcal F}}
\def\P{{\mathcal P}}

\def\newpage{\vfill\eject}
\title{The line graph of a bipartite graph has its edge set partitioned into two comparability graphs}
\author{András Gyárfás, Márton Marits, Géza Tóth}
\date{arxiv:2408.13523v2; CC BY 4.0}
\begin{document}
\maketitle
\noindent\textit{Editorial note.} The counted unit is the article's own Theorem bipline at source lines 259-260 together with the author's proof that opens with the bold header Proof. at source line 265 and closes with the square at source line 268; the author exceeds the usual placement by inserting a medskip at line 263 between the statement and the proof, so the pairing is by adjacency through blank lines and that medskip only.\par\medskip
\noindent\textit{Editorial note.} The retained context at source lines 117-120 and 128-129 is the article's own description of comparability graphs and its definitions of p(G) and c(G); the retained context at source lines 255-257 is its own definition of a line graph and its notation LBIP for the line graphs of bipartite graphs.\par\medskip
\noindent\textit{Editorial note.} The source preamble cannot be retained as it stands: it loads rawfonts, psfrag, verbatim, tikz with several libraries, hyperref with the pagebackref option and a dvi driver option for graphicx, none of which the retained body spans use, and it redefines backref commands that only exist because of that hyperref option. The retained preamble file therefore carries, line for line, source line 2, source line 3, source line 4, source line 16, source line 17, the theorem declarations of source lines 20-30 and the author macro definitions of source lines 32-41; the only textual change anywhere is the removal of the dvipdfm driver option from the graphicx line, because the assembly forces the xelatex engine. No new macro, package or declaration was introduced.\par\medskip
\noindent\textit{Editorial note.} No retained line contains a citation or a cross-reference, so the material carries no bibliography and needs no external navigation pointer.\par\medskip
\noindent\textit{Editorial note.} The article's Theorem twocover, Theorem promelsteger, Theorem gensplit, Theorem unit, Theorem co-triangulated, Theorem loglog, Theorem ub, Lemma rec, the corollaries genupbound, shift, doblesh and p', and Proposition H9 are not counted here.\par\medskip
\noindent\textit{Allowed standard prerequisites (editorial).}\begin{itemize}
\item Comparability graph and the partition number p(G). A comparability graph has an edge between two vertices exactly when the corresponding elements of a partially ordered set are comparable, and p(G) is the minimum number of edge-disjoint comparability subgraphs of G whose union is the whole edge set; retained at source lines 117-120 and 128.. The counted theorem asserts p(G) <= 2 for a class of graphs, so the class used in the partition and the parameter bounded must be available.
\item Line graphs and the class LBIP. The line graph of a graph G has the edges of G as its vertices, two of them adjacent when the corresponding edges of G share a vertex, and LBIP denotes the class of line graphs of bipartite graphs, which the article records at source lines 255-257 as a well-known class of perfect graphs.. Identifies the graphs of the counted theorem as line graphs of bipartite graphs and supplies the bipartition [A,B] of the underlying bipartite graph that the proof uses to label the edges.
\end{itemize}\medskip
Here we focus on a well known family of perfect graphs, the {\em comparability graphs},
which can be obtained from a partially ordered set as follows.
Vertices correspond to the elements of the partially ordered set,
and two vertices are connected by an undirected edge if and only if the corresponding elements

Let $p(G)$ denote the minimum number $m$ such that $E(G)=\cup_{i=1}^m E(G_i)$ where the $G_i$-s are edge-disjoint comparability subgraphs of $G$.
Dropping  the condition that the $G_i$-s are edge-disjoint, we define $c(G)$ as the smallest number of comparability subgraphs of $G$ needed to {\em cover} $E(G)$.

The {\em line graph} of a graph $G$ is obtained by representing the edges of $G$ by vertices and
considering two such vertices adjacent if the corresponding edges of $G$ have a common vertex.
The line graphs of bipartite graphs, LBIP, form a well-known class of perfect graphs.

\begin{theorem}\label{bipline} For $G\in $LBIP, $p(G)\le 2$.
\end{theorem}

\noindent{\bf Proof. }  Let $G$ be the line graph of a bipartite graph $H=[A,B]$.
The edges of $G$ are labelled with $a$ ($b$) if the the corresponding edges of $H$ intersect in $A$ (in $B$).
Clearly, the edges labelled with $a$ (with $b$) are unions of complete graphs with no edges between them,
thus they give the required partition of $E(G)$ into two comparability graphs. \qed

\end{document}
